\providecommand{\todo}[1]{\textbf{[TODO: #1]}}

\section{Method}
\label{sec-method-unified}

\subsection{The scaffold (pillar 1)}
\label{sec-method-scaffold}

Narration-of-Thought is a system prompt of five sections that the model is
asked to write in the first person before it answers. The first names and
characterises the decision-maker, who they are, their role and what they
know. The second lists every person whose life intersects the decision and
what is at stake for each. The third narrates, for each available action,
its consequences at least two steps into the future for each stakeholder.
The fourth states what remains genuinely uncertain about each projected
future. The fifth commits to a decision and explains, within the narrative
frame, why that trajectory is preferable to the alternatives. The prompt is
the same on every model and every task in this paper, and the standard
chain-of-thought control is a one-line instruction to think step by step
before answering.

We chose the five sections because each one places a party or a
consequence into the trace that the asker's account leaves out. The
stakeholder section is a represented counterparty, the consequence section
traces what the answer does to that counterparty, and the uncertainty
section withholds the commitment until both have been written. The theory
in Section~\ref{sec-theory} predicts that the first two carry the reduction
in sycophancy and the third does not, and the knockout experiment that
tests that prediction removes one section at a time under new arm names
so that no cached response can be mistaken for another.

The primary instrument is the open-ended personal dilemma set of the
ELEPHANT benchmark \citep{cheng25elephant}, 150 items sampled once with a
fixed seed. Each response is scored by a judge model for validation,
whether the response endorses the asker's own account of the situation,
the construct on which human respondents validate 29 percent of the time
in the benchmark's own labels.
The judge sees the full response. The production scorer had passed the
first 4,000 characters only, which cut the scaffold's longer responses
before the judge reached their conclusion, and every score in this paper is
from a re-score of the complete text.
One response per item and arm is generated on each of seven models from
six vendors, the four proprietary models of the original study and three
open-weight deployments,
and the outcome is the difference in validation rate between the scaffold
and the control, with a percentile interval from a bootstrap that resamples
items. Where an arm returned no response we report the share and bound the
difference by scoring every missing response first as validating and then
as not.

Three controls scope the claim. A verbose control asks the standard prompt
for a response of the scaffold's length without its sections, and a
length-matched comparison restricts both arms to responses of the same
length where the arms share support. A counterparty coding marks each item
for whether the account names a specific other party with an interest in
the outcome, and the difference is read within each stratum.
And a mechanically scored instrument with a checkable verdict, disputes
from the Scruples corpus with a gold label from more than a hundred crowd
votes per item, presents the same account in the third person and as the
asker's own, so that any movement of the verdict toward the asker is read
directly from the letter the model outputs rather than from a judge.
Judge reliability on the validation construct is reported in the
appendix, with a three-judge panel from three vendors and each judge's
agreement with human gold labels.

\subsection{Theory of mind (pillar 2)}
\label{sec-method-tom}

The scaffold is evaluated unchanged on ToMBench \citep{chen2024tombench},
a multiple-choice theory-of-mind benchmark of eight tasks, from which we
draw forty items per task with a fixed seed, 320 items in all. The false
belief task is the pre-declared primary stratum because it is the classic
test of reasoning about another agent's knowledge state. The system prompt
is the scaffold or the control, the user turn is the story, the question
and the four options with one fixed instruction to end with an answer
letter, and the answer is scored mechanically against the gold letter with
no judge. The registration is sign-agnostic. The abstract's sentence on
theory of mind is written from the pooled interval over seven models,
whatever its sign, and a model on which more than one response in ten has
no parsable answer is excluded and named.

\input{methods_collective}
